\chapter{CONCLUSÃO}
\label{cap:conclusao}

Este trabalho investigou em que medida um modelo de linguagem de pequeno porte, executado localmente e restrito a escolher entre candidatos recuperados da tabela oficial, consegue apontar classificações NCM inválidas ou incorretas no cadastro de produtos de um ERP e sugerir a classificação correta, em tempo compatível com o uso diário e sob supervisão humana. Para responder a essa pergunta, foi desenvolvido e avaliado um protótipo integrado ao ERP Lux Auto, executado inteiramente em um computador convencional, sem custo de API e sem alterações no código ou na estrutura do ERP.

A resposta obtida é positiva para a detecção de códigos inválidos e parcial para a detecção de códigos incorretos e para a sugestão. A detecção de códigos inválidos foi resolvida de forma determinística, sem modelo de linguagem: todos os códigos vazios, malformados, inexistentes e extintos dos casos de avaliação foram apontados, sem nenhum falso positivo entre mais de 41 mil mercadorias reais. A detecção de códigos válidos, porém incorretos, que depende do modelo, teve acurácia de 86,6\% no desenvolvimento e de 88,8\% nos erros reais de uma loja, na qual apontou todos os 40 produtos cadastrados com um código ``coringa''; no conjunto de teste, a acurácia foi de 86,0\%, com 6,0\% de alarmes falsos, mas apenas 78,0\% dos erros plantados foram apontados, pois erros entre itens vizinhos de uma mesma posição continuam difíceis de perceber. Na sugestão, a versão final acertou o código em primeiro lugar em 48,0\% dos produtos de teste e entre as três opções em 58,0\%, em cerca de 14 segundos por produto, abaixo da meta de 75\% no conjunto completo, mas com 83,3\% de acerto em primeiro lugar nos produtos automotivos, foco do ERP. Em todos os casos, nenhuma sugestão continha código inexistente, e nenhum código foi gravado sem aprovação humana.

Os objetivos específicos foram atendidos da seguinte forma: as inconsistências fiscais foram delimitadas a partir da experiência no suporte técnico, e a NCM foi escolhida como foco; a tabela NCM foi estruturada em uma base local com busca textual; o protótipo foi implementado com detecção de códigos inválidos e de códigos possivelmente incorretos, sugestão, execução em segundo plano, notificação e aprovação humana; foram construídos conjuntos separados de desenvolvimento e de teste, incluindo códigos inválidos de diferentes tipos e erros reais de classificação; e o protótipo foi avaliado quanto à detecção, ao acerto e ao tempo, com comparação entre três modelos, entre arquiteturas e contra linhas de base sem modelo de linguagem.

As principais contribuições do trabalho são: uma arquitetura de busca seguida de escolha restrita, na qual o modelo de linguagem não acessa o banco nem pode responder com códigos inexistentes; a medição da contribuição do modelo, que elevou o acerto em primeiro lugar de 28,0\% para 48,0\% em relação à melhor combinação de busca e base de conhecimento sem ele; o uso de produtos já classificados por outras empresas como base de conhecimento, sem treinamento, que elevou substancialmente o acerto no setor automotivo; um método de verificação de códigos válidos que encontrou um erro real e sistemático no cadastro de uma loja; e um protocolo de avaliação com conjunto de teste congelado e medido uma única vez, que mostrou o quanto a avaliação nos próprios casos de ajuste superestima o desempenho.

\section{Trabalhos Futuros}

A partir dos resultados e das limitações observadas, sugerem-se as seguintes continuidades:

\begin{itemize}
    \item \textbf{Demais categorias fiscais}: estender a arquitetura ao CST, ao CFOP e à relação entre Inscrição Estadual e Indicador de Contribuinte, cuja detecção pode, em boa parte, seguir regras determinísticas semelhantes às da NCM;
    \item \textbf{Bases de outros setores}: ampliar a base de conhecimento com empresas de outros ramos, que hoje concentram a maior parte dos erros, e medir o ganho em cada setor;
    \item \textbf{Avaliação com gabarito profissional}: medir o sistema em produtos reais de clientes com o gabarito conferido por um profissional de classificação fiscal, o que permitiria também conferir os alarmes da detecção de NCM incorreto;
    \item \textbf{Vocabulário técnico e notas explicativas}: construir, a partir das aprovações dos operadores, um dicionário de correspondências entre nomes comerciais e termos da NCM, e incorporar à busca as Notas Explicativas do Sistema Harmonizado, atacando os erros de capítulo, que são a maior parte dos erros restantes;
    \item \textbf{Abreviações aprendidas}: ampliar a tabela de abreviações com as encontradas nas bases reais, que causaram parte dos alarmes falsos da detecção;
    \item \textbf{Modelos maiores ou ajustados}: avaliar modelos maiores, quando o hardware permitir, o modo de raciocínio explícito do Qwen3 ou um modelo pequeno ajustado com exemplos de classificação, comparando ganho de acerto e custo de execução;
    \item \textbf{Integração com o ERP}: substituir a consulta periódica por um aviso no próprio ERP no momento do cadastro, caso a desenvolvedora opte por essa integração.
\end{itemize}
